\documentclass[11pt]{article}
\usepackage[margin=0.9in]{geometry}
\usepackage{amsmath,amssymb,amsthm,booktabs,array}
\usepackage[T1]{fontenc}
\usepackage{lmodern,microtype}
\usepackage[colorlinks=true,linkcolor=blue,urlcolor=blue]{hyperref}
\newtheorem{theorem}{Theorem}[section]
\newtheorem{lemma}[theorem]{Lemma}
\newtheorem{proposition}[theorem]{Proposition}
\newtheorem{corollary}[theorem]{Corollary}
\theoremstyle{remark}
\newtheorem*{scope}{Scope}
\newcommand{\tr}{\operatorname{Tr}}
\newcommand{\norm}[1]{\left\|#1\right\|}
\newcommand{\Z}{\mathbb Z}
\newcommand{\R}{\mathbb R}
\newcommand{\E}{\mathbb E}
\newcommand{\supp}{\operatorname{supp}}
\title{Toward a four-dimensional Yang--Mills construction\\
\large Notebook 02: preserving diagonal reflection positivity}
\author{Research working document}
\date{16 October 2026}
\setlength{\emergencystretch}{5em}
\begin{document}
\maketitle
\begin{abstract}
We give a fixed-coupling argument transferring diagonal reflection
positivity to the infinite-volume limits of the original coordinate
Wilson tori of side $2L^M$, with $L$ odd. The construction uses exact
Wilson slab gluing, alternating coordinate site/link cuts, local
gauge-invariant layer removal, and convex separation. It requires
neither uniqueness of the infinite-volume state nor estimates uniform
as the coupling tends to infinity. This settles one regulator-level
interface in the proposed continuum construction, subject to scrutiny
of the proof below. It does not establish continuum rotations,
renormalized source bounds, nontriviality, or a mass gap.
\end{abstract}

\noindent\textbf{Evidence and status.}
This is an explicit research proof, internally checked in separate
passes for the local gauge argument, geometry, slab normalization, and
final limit argument; it has not received external mathematical review.
The comparison architecture is adapted from the QED3 material in the
\href{https://github.com/ethanwsussman/Yang--Mills-LLM-experiment/tree/e392ef45013a581664bab76b90867170b30fe10e}
{pinned repository}, rather than claimed as a new general method.
The Wilson and odd-$L$ details are supplied below. Notebook 01 retains
the other construction obligations.

\section{Precise statement and order of limits}

On independent positively oriented links of $\Z^4$ put $U_e\in SU(2)$,
with $U_{\bar e}=U_e^{-1}$. Haar measures have mass one. At fixed
$0<\beta<\infty$, write
\[
 k_\beta(g)=\exp\{(\beta/2)\tr g\},\qquad
 Z_\Gamma[F]=\int F(U)\prod_{p}k_\beta(U_p)\prod_e dU_e.
 \tag{1.1}
\]
Every unoriented elementary plaquette occurs once. Let
$C_n=\Z^4/n\Z^4$, $n=2L^M$, for a fixed odd integer $L\ge3$.
The normalized Wilson measure is denoted $\mu_{C_n}$.
Small tori with degenerate plaquette identifications can be discarded.

A lattice reflection acts on oriented links and hence on configurations.
For its ordinary pullback $rF=F\circ r$, put
$\Theta F=\overline{rF}$. Reflection positivity (RP) means
\[
 \nu((\Theta F)F)\ge0
 \tag{1.2}
\]
for every continuous local $F$ supported in the closed positive half.
Equivalent finite sums follow by replacing $F$ with their sum.

\begin{theorem}\label{thm:main}
Every locally weak subsequential limit $\nu$ of $\mu_{C_n}$ is
reflection invariant and reflection positive across
$x_1-x_2=0$. By its lattice translations and signed coordinate
permutations, the same holds at every lattice-preserving diagonal
mirror $x_i\pm x_j=c$, $i\ne j$, $c\in\Z$.
The original coordinate RP properties are retained as well.
\end{theorem}

The proof constructs, for each separating observable $O$, auxiliary
tilted measures. Its limits are strictly ordered:
\[
 \text{fix }\beta,O,\text{ and all cell sizes};\quad
 n\longrightarrow\infty;\quad t\downarrow0;
 \quad\text{then approximate local observables.}
 \tag{1.3}
\]
No interchange with $\beta\to\infty$ is made. Applying the theorem
at each lattice spacing is compatible with a volume-first continuum
construction. Arbitrary simultaneous volume/cutoff limits would need
additional estimates.

\section{Exact Wilson slabs and reflected cell inequalities}

\subsection{Diagonal boundary variables}
Take height $s=x_1-x_2$ and an integer mirror $s=c$.
Links in directions 1 and 2 end on this mirror or avoid it; their
interiors never cross it. Links in directions 3 and 4 can lie in it.
A $(1,2)$ plaquette based at height $b$ has heights
$b,b+1,b-1,b$ and crosses the mirror exactly when $b=c$.
It splits into two length-two paths between the same boundary vertices.
Plaquettes in directions $(1,\mu)$ or $(2,\mu)$, $\mu=3,4$,
have height range one and lie on one side. A $(3,4)$ plaquette
can lie wholly in the mirror. For $s=x_1+x_2$ the crossing
$(1,2)$ plaquette is centered at $c=b+1$; the same classification holds.

For completeness the following feature construction works for $SU(N)$.
Let $u(A)$ be the real vector of real and imaginary matrix entries, so
$u(A)\cdot u(B)=\operatorname{Re}\tr(AB^\dagger)$ and
$|u(A)|^2=N$. Set $c_\beta=\beta/N$ and
\[
 \Phi_\alpha(A)=
 \left(\frac{c_\beta^{|\alpha|}}{\alpha!}\right)^{1/2}
 u(A)^\alpha,\qquad \alpha\in\mathbb N^{2N^2}.
\]
Then
\[
 \langle\Phi(B),\Phi(A)\rangle
 =e^{c_\beta\operatorname{Re}\tr(AB^\dagger)},
 \qquad\norm{\Phi(A)}^2=e^\beta.
 \tag{2.1}
\]
The coefficients are real and the series converges absolutely by
Cauchy--Schwarz. If $A_+,A_-$ are the two half-path holonomies with
common initial and final vertices, their feature contraction is exactly
the crossing plaquette weight.

Work on a cylinder with finite tangential period lattice
$\Gamma_\parallel\subset\{s=0\}$. At $s=c$ take boundary Hilbert space
\[
 \mathcal H_c=L^2(SU(2)^{E_c},db)
       \widehat\otimes\bigotimes_{p\in P_c^\times}\mathcal F_p,
 \tag{2.2}
\]
where $E_c$ are links in the sheet, $P_c^\times$ the crossing
plaquettes, and $\mathcal F$ the complexified feature space above.
Give each neighboring slab a square root of the product $W_c(b)$
of weights of plaquettes contained in the sheet.

For a bounded function $F$ strictly inside a slab $[c_0,c_1]$,
define $S[F]:\mathcal H_{c_0}\to\mathcal H_{c_1}$ by integrating
its interior links, multiplying every one-sided plaquette weight,
the two boundary halfweights, and the feature coefficient for each
inside half-path at either boundary. More explicitly, its kernel is
\[
 \sqrt{W_{c_1}(b_1)W_{c_0}(b_0)}
 \int F\prod_{p\text{ assigned to slab}}k_\beta(U_p)
       \prod_{p\text{ at boundaries}}\Phi_{\alpha_p}(A_{p,\mathrm{in}})
       \prod_{e\text{ in slab}}dU_e.
 \tag{2.3}
\]
Widths are at least four, as they are in every application here.
The constant feature norm and bounded finite-volume Wilson weights,
followed by Minkowski's inequality, prove that (2.3) is a
Hilbert--Schmidt kernel, also for complex $F$.

\begin{lemma}[Gluing and adjoints]\label{lem:glue}
Adjacent slab kernels compose to the kernel of their union with the
product of their integrands. Reflection of a slab takes its operator
adjoint, with the natural boundary translation identification.
No extra normalization factors occur.
\end{lemma}
\begin{proof}
At a common sheet, kernel composition integrates each boundary link
once. Feature contraction restores each crossing plaquette by (2.1);
the two halfweights multiply to $W_c$. Every other link or plaquette
already belongs to exactly one slab. All Haar measures have mass one.
Reflection exchanges the two half-paths and kernel arguments;
the real feature coefficients give the complex adjoint.
If the width is $b$, reflection in its upper mirror maps its lower
boundary to the translate by $q=bv$, where $v=e_1-e_2$ or $e_1+e_2$.
Thus, with the translation unitary $U_q$,
$S[\Theta F]=U_qS[F]^*$. This uses an integer translation even
when $b$ is odd; the two original boundary spaces need not be
identified by a normal translation of length $b/2$.
\end{proof}

If $\Gamma=\Gamma_\parallel\oplus\Z\gamma$,
$s(\gamma)=p$, the same graph gluing gives
\[
 Z_\Gamma[F_1\cdots F_J]
 =\tr(U_\gamma^{-1}S[F_J]\cdots S[F_1]).
 \tag{2.4}
\]
The closing translation may have a tangential component. It is a
unitary relabeling, not a new action at the seam.

\subsection{Even folds and four successive directions}
Suppose $p=Jb$ with $J$ even. The straightened period is
$\gamma_0=(p/2)v\in\Z^4$. Denote by $\operatorname{Fold}_sF$
the product of successive antilinear reflected copies of one slab
function, on the straightened torus $\Gamma_0$.
Each reflected pair contributes a translate of $S[F]^*S[F]$.
Their translation unitaries telescope to $U_{\gamma_0}$ and cancel
the closing unitary in (2.4). Hence
\[
 Z_{\Gamma_0}[\operatorname{Fold}_sF]=\norm{S[F]}_J^J\ge0,
 \qquad
 |Z_\Gamma[\textstyle\prod_jF_j]|
 \le\prod_j Z_{\Gamma_0}[\operatorname{Fold}_sF_j]^{1/J}.
 \tag{2.5}
\]
The second inequality is Schatten H\"older. Different boundary
Hilbert spaces are allowed; the operators can be viewed as cyclic
blocks in their orthogonal direct sum.

Coordinate versions, including unlike site/link boundaries, are
proved in Section~\ref{sec:coordinate}. For four orthogonal families
with even successive counts $J_1,\ldots,J_4$, apply (2.5) one direction
at a time. Reflections in one direction preserve the other heights.
After each completed gluing one has an ordinary Wilson action, so
there is no simultaneous assignment of corner halfweights. The
terminal factor is the full reflected array of one original cell.
Writing $N_{\rm cell}=\prod_iJ_i$, this proves
\[
 \left|Z_\Gamma[\textstyle\prod_cF_c]\right|
 \le\prod_c
 Z_{\Gamma_{\rm final}}[\operatorname{Fold}_{\rm cells}F_c]^{1/N_{\rm cell}}.
 \tag{2.6}
\]
Its terminal factors are nonnegative by the last use of (2.5).
All cell integrands used below have support strictly inside their cells.

The same boundary factorization with two reflected halves proves RP
on a straight diagonal torus, including bounded positive source
weights that split into reflected half products. Reflection invariance
follows directly from its graph symmetry.

\section{Coordinate physical operators and the insertion bound}
\label{sec:coordinate}

On a finite spatial three-torus with edge set $E$, let
$\mathcal H=L^2(SU(2)^E)$ and let $P$ be Haar averaging over the
spatial endpoint gauge group. Write $\mathcal H_{\rm phys}=P\mathcal H$.
Let $C=\bigotimes_{e\in E}C_\beta$ be electric convolution with
$k_\beta(vu^{-1})$, and let $V$ multiply by all spatial plaquette
weights. They commute with $P$. The positive convolution eigenvalues
are $2I_{2j+1}(\beta)/\beta$; their rapid decay makes $C$ trace class.

On the physical space define
\[
 B=V^{1/2}C^{1/2},\quad A_l=B^*B=C^{1/2}VC^{1/2},
 \quad A_s=BB^*=V^{1/2}CV^{1/2}.
\]
Temporal-link Haar integration gives $CP$ at every transition:
its plaquette is $\tr(U_ea_yU_e'{}^{-1}a_x^{-1})$.
Consequently the ordinary coordinate partition function is exactly
\[
 Z_\Gamma(N)=\tr_{\mathcal H}(CPV)^N
           =\tr_{\rm phys}A_l^N=\tr_{\rm phys}A_s^N.
 \tag{3.1}
\]
There is no gauge-volume factor. A link cut has a factor $C^{1/2}$
on either side; a site cut splits $V$ into $V^{1/2}$. Halfsteps
are $B$ and $B^*$. Thus a slab can have unlike boundary types.

\begin{lemma}[Bounded sources in coordinate slabs]\label{lem:schatten}
Let a slab have length $\ell\in\tfrac12\mathbb N$ and bounded real
gauge-invariant source $R$, $\norm R_\infty\le M$, whose support crosses neither cut.
For every even integer $p\ge2$, its boundary operator satisfies
\[
 \norm{S_{R,t}}_p\le e^{|t|M}Z_\Gamma(p\ell)^{1/p}.
 \tag{3.2}
\]
If $J$ is even, $\sum_j\ell_j=N$, and bounded operators $X_j$
are placed at the intervening cuts, the unnormalized insertion obeys
\[
 \left|\tr(X_JS_{R_J,t}\cdots X_1S_{R_1,t})\right|
 \le e^{|t|\sum_jM_j}\prod_j\norm{X_j}
       \prod_j Z_\Gamma(J\ell_j)^{1/J}.
 \tag{3.3}
\]
\end{lemma}
\begin{proof}
Glue $p$ alternating copies of the slab and its adjoint.
At link joins the factors give $C^{1/2}PC^{1/2}=CP$;
at site joins they give $V^{1/2}PV^{1/2}=VP$, with that
projector absorbed by adjacent physical temporal cores.
The trace of $(S^*S)^{p/2}$ is the ordinary Wilson integral
on length $p\ell$ with $p$ bounded source copies. Pointwise
bounding their product gives (3.2). Finite representation truncations
or the bounded glued integral justify the kernel limit.
No pointwise positivity of the kernel of $C^{1/2}$ is asserted.
Schatten H\"older gives (3.3).
\end{proof}

For general bounded complex gauge-invariant slab integrands the
operator is defined linearly: split real and imaginary parts, and
write a bounded real function as a difference of two strictly positive
bounded functions, each an exponential of a bounded real function.
The kernels glue linearly and reflection is the adjoint. This proves
the coordinate version of (2.5) for half-integer widths and even counts.
A tangential closing shift acts unitarily on both $\mathcal H$ and
$\mathcal H_{\rm phys}$ and gives the same twisted-trace formula.

For a selected finite set of spatial links $S$ at a link cut, replace
the electric kernel $K_S$ by that of $C_S^{1/2}XC_S^{1/2}$ and divide
by $K_S>0$. Include temporal parallel transports at its endpoints.
The resulting local insertion is denoted $G_X$; $G_I=1$.
Throughout the rest of this note, $X$ is restricted to finite
Peter--Weyl kernel operators and their unitization that commute with
endpoint gauge transformations. Their insertion kernels and $G_X$
are continuous, $X\otimes I$ commutes with $P$, and $G_X$ is
gauge invariant. This class contains the self-adjoint combinations and
$I+isA$ used below. Disjoint spatial supports at a common cut give an
operator norm at most the product of the individual norms.
Lattice reflections and relabelings replace $X$ by adjoints or
unitary/antiunitary conjugates, preserving its norm.

\section{Local gauge-invariant reduction to two slices}

\begin{lemma}[Finite derivative loss]\label{lem:convolution}
For a finite Peter--Weyl polynomial $f$ on $SU(2)$ of maximum spin
$r$, $T_f=C_\beta^{-1}M_fC_\beta$ maps $C^\infty$ continuously to
itself and
\[
 \norm{T_f\psi}_{H^s}\le C_{\beta,f,s}\norm\psi_{H^{s+2r}}.
 \tag{4.1}
\]
The tensor version holds on any fixed finite edge set.
\end{lemma}
\begin{proof}
The spin-$j$ convolution eigenvalue is
$c_j=2I_{2j+1}(\beta)/\beta>0$. Multiplication by $f$ connects
only $|k-j|\le r$. In
\[
 I_m(\beta)=\frac{(\beta/2)^m}{m!}S_m,\qquad
 1\le S_m=\sum_{l\ge0}
 \frac{(\beta^2/4)^l m!}{l!(m+l)!}\le e^{\beta^2/4},
\]
ratios on that finite band grow at most as $C(1+j)^{2r}$.
Each multiplication block has norm at most $\norm f_\infty$;
weighted square summation proves (4.1). Sobolev embedding and the
finite tensor product give the stated smooth mapping property.
\end{proof}

We record the gauge step explicitly because replacing the law by a
gauge-fixed one would obscure the support issue. Suppose a local
gauge-invariant polynomial $P$ has lowest coordinate height $j$.
For its bottom spatial edges $e=(x,y)$ put
\[
 b_e=U_{(j,x),(j,y)},\quad a_x=U_{(j,x),(j+1,x)},\quad
 u_e=a_x^{-1}b_ea_y,\quad v_e=U_{(j+1,x),(j+1,y)}.
\]
Include auxiliary $a_x$ at all necessary endpoints. Gauge transformation
$g_{j,x}=a_x^{-1}$ sets those temporal links to one and leaves every
upper link unchanged. It proves the identity of observables
\[
 P(U)=\widetilde P(u,U_{\ge j+1})
      =\sum_i f_i(u)P_i(v,\eta).
 \tag{4.2}
\]
There is no change of the probability measure.

Condition on all links except the finite collection $v$.
Its downward plaquettes have trace $\tr(u_ev_e^{-1})$.
Every other plaquette containing $v$ is spatial at $j+1$ or goes
up to $j+2$. Include each such complete plaquette once in $W(v;\eta)$.
The density factors as
\[
 H(\mathrm{exterior})K_S(v,u)W(v;\eta).
\]
$W$, its inverse, and all required derivatives are bounded uniformly
in ambient volume at fixed $\beta$ and fixed support.
Since $C_ST_f=M_fC_S$ and $K_S(v,u)=K_S(u,v)$, define
\[
 Q(v,\eta)=W(v;\eta)^{-1}
          \sum_i T_{f_i}[W(\cdot;\eta)P_i(\cdot,\eta)](v).
 \tag{4.3}
\]
Conditional integration proves $\mu(P)=\mu(Q)$ exactly.
$Q$ is smooth and reads only links at heights at least $j+1$,
reaching up at most to $\max(j+2,\text{old top})$.
It is gauge invariant: $u_e$ and $v_e$ transform by the same
upper endpoint gauges, $C_S$ commutes with their action, and
$T_{hf}(hB)=h(T_fB)$. Thus (4.3) descends directly to the original
upper links without reintroducing the removed links.

Now tilt by $e^{tV_n}$, where $V_n$ is a fixed-spacing array of
reflected copies of a smooth bounded local source. Only a uniformly
bounded number of copies depend on the integrated $v$; call their sum
$R_n$. It may also read lower exterior links. With those frozen,
the exact tilted replacement is
\[
 Q_{n,t}=W^{-1}e^{-tR_n}
           \sum_iT_{f_i}[e^{tR_n}WP_i].
 \tag{4.4}
\]
Finite derivative loss and compactness give
$\norm{Q_{n,t}-Q}_\infty\le C_Pt$ for $0\le t\le1$,
uniformly in $n$. The analogous top-layer removal uses
$u_e=a_xb_ea_y^{-1}$ and goes down one slice.

\begin{proposition}[Paired two-layer reduction]\label{prop:layers}
For every continuous gauge-invariant local $P$, every integer $a$,
and $\epsilon>0$, there is a finite local gauge-invariant polynomial
$Y_\epsilon$ supported in $[a,a+1]$, and $C_\epsilon<\infty$, with
\[
 \begin{split}
 |\mu_{n,t}(P)-\mu_{n,t}(Y_\epsilon)|&\le\epsilon+C_\epsilon t,\\
 |\mu_{n,t}(rP)-\mu_{n,t}(rY_\epsilon)|&\le\epsilon+C_\epsilon t,
 \end{split}
 \tag{4.5}
\]
where $r$ is reflection in $a+1/2$. No invariance of the tilted
measure or source array under $r$ is required. The estimates hold
for all sufficiently large $n$, once the finite neighborhoods produced
by reduction and their reflections embed without identifications.
\end{proposition}
\begin{proof}
First uniformly approximate a continuous input by a finite Peter--Weyl
polynomial and average over its finite endpoint gauge group, charging
this initial approximation to the total $\epsilon$ allowance.
Remove bottom layers up to $a$, then top layers down to $a+1$.
The first operations never exceed the retained top bound since
$j+2\le a+1$; the second never drop below $a$ since $j-2\ge a$.
Each operation enlarges spatial support by finitely many plaquettes.
After each step approximate its smooth output uniformly by a finite
Peter--Weyl polynomial and Haar-average over the finite endpoint gauge
group. The average preserves both support and finite representation
content. Choose the finite list of approximation errors with total
below $\epsilon$. Expectation errors telescope, so subsequent operator
norms do not amplify earlier errors. The finite sum of tilted errors
is $C_\epsilon t$. Reflect the same unperturbed construction for
$rP$; its tilted error proof needs only the same local derivative
bounds, not reflection symmetry of the tilt. This proves (4.5).
\end{proof}

\section{The two frames on the original odd-$L$ family}
\label{sec:geometry}

Put $a=n/2=L^M$, $e_-=e_1-e_2$, $e_+=e_1+e_2$, and introduce
\[
 \Gamma_D=\langle ae_-,2ae_+,ne_3,ne_4\rangle,\qquad
 \Gamma_T=\langle ae_1,4ae_2,ne_3,ne_4\rangle.
 \tag{5.1}
\]
Both covolumes are $n^4$. These are auxiliary tori; $C_n$ is unchanged.

Let $O$ be a real gauge-invariant local Peter--Weyl polynomial.
Choose a source radius $R_O$ large enough to contain its full link
support and let $h=L^u>16R_O+16$. Fine diagonal cells have widths $h$
in $(x_1-x_2,x_1+x_2,x_3,x_4)$, with integral diagonal offsets
and half-integral transverse offsets. The source is strictly inside
one cell; $V_n$ is the sum of its reflected copies.

For an insertion box $B=y+[-r,r]^4$, choose fixed powers
$q=L^v$, $b=qh$, $H=L^z$, with $b\mid H$, and put $w=b/2$, so
\[
 w\ge2r+6h,\qquad H/w\ge16.
 \tag{5.2}
\]
For large $M$, $h\mid b\mid H\mid a$.
Both $b,H$ are odd; $w$ is strictly half-integral. Coordinate
walls spaced by $w$ therefore alternate between site and link cuts.

\subsection{Periods and counts}
On $C_n$ the first diagonal height has period $n$ and tangential
cross-section generated by $ne_+,ne_3,ne_4$. Straightening replaces
the closing vector $ne_1$ by $ae_-$. The subsequent height periods
are $2n,n,n$. Thus the first fold ends on $D_n$ with even counts
\[
 (n/h,2n/h,n/h,n/h),\qquad N=2n^4/h^4.
 \tag{5.3}
\]
The same fold on $D_n$ with width $H$ has
$N'=2n^4/H^4$ cells. The diagonal height transformation has
determinant two, explaining the cell volume $h^4/2$.
Both original and adapted periods belong to the wall group's
translation subgroup generated by $he_-,he_+,2he_3,2he_4$;
all reflected arrays are therefore periodic.

For the first coordinate height on $D_n$, the period is $a$, with
cross-section $\langle4ae_2,ne_3,ne_4\rangle$ and closing vector
$ae_1-ae_2$. Its tangential shift is the unitary relabeling $-ae_2$.
Straightening gives $T_n$. With $A=a/b$, the coordinate counts are
\[
 (2A,8A,4A,4A),\qquad N''=256A^4=n^4/w^4,
 \tag{5.4}
\]
all even. In units of $w$ the original period generators are
$(2A,-2A),(4A,4A),4Ae_3,4Ae_4$. First reduce the first index by
the first generator, then the second by $(0,8A)$, then the transverse
ones. This proves that the brick (5.4) represents both $D_n$ and
$T_n$, including the intermediate cellwise folding geometry.

\subsection{Clean planes and distinct insertion cells}
The reflected source centers have at most four residues modulo $h$
in either first coordinate, and at most two in the transverse ones.
Indeed, writing $A_\pm=x_1\pm x_2-c_\pm$, their first coordinates
are
\[
 x_1'=x_1-A_-\epsilon_--A_+\epsilon_++h(z_-+z_+),\quad
 x_2'=x_2+A_-\epsilon_--A_+\epsilon_++h(z_+-z_-).
 \tag{5.5}
\]
For two adjacent half-integer planes $c,c+1$, one source-center
residue forbids at most $2R_O+2$ of the $h$ half-integer classes
modulo $h$. Hence two adjacent clean link planes exist in each
coordinate direction. Here clean means that neither plane meets
any source box.

Choose coarse offsets from the fine wall families so that
\[
 \begin{aligned}
 y_1-y_2-c'_-&\in[2w,2w+h),&
 y_1+y_2-c'_+&\in[4w,4w+h),\\
 y_\mu-c'_\mu&\in[w,w+h),&&\mu=3,4.
 \end{aligned}
\]
For the resulting coarse wall group, put $m=H/w\ge16$, an even
integer. Every image center has form
\[
 gy=y+w\kappa(g)+s(\epsilon),\qquad |s(\epsilon)|_\infty<2h,
\]
where
\[
 \begin{split}
 \kappa(g)={}&(-2\epsilon_--4\epsilon_+,
               2\epsilon_--4\epsilon_+,-2\epsilon_3,-2\epsilon_4)\\
 &+m(z_-+z_+,z_+-z_-,2z_3,2z_4).
 \end{split}
 \tag{5.6}
\]
The sixteen parity vectors are distinct; coordinate differences
are at most six. Since $m\ge16$, equality of indices forces equality
of the parity vectors and then of every $z$. Thus $g\mapsto\kappa(g)$
is injective.

Coordinate walls of width $w$ have residues $o,o+h/2$ modulo $h$.
Choose $o$ modulo $h/2$ on the half-integer mesh. A source residue
excludes at most $4R_O+2$ classes; all four exclude at most
$16R_O+8<h$. An allowed representative can be selected in
\[
 (y_\mu+r+2h-w,\ y_\mu-r-2h),
\]
whose length is at least $2h$. Equations (5.2) and (5.6) then give
\[
 gB\subset o+w\kappa(g)+(0,w)^4.
 \tag{5.7}
\]
Every source box is also strictly inside one coordinate cell.
The periods in $\Gamma_D$ belong to the coarse diagonal translation
group because $H\mid a$. If two indices in (5.6) differ by a
period, translate one group element by that period and use
injectivity. They represent the same coarse orbit. Consequently
exactly $N'$ distinct coordinate cells contain amplified insertions.

After the coordinate fold of one occupied cell, all insertion cuts
have the same coordinate direction. If its relative cut position is
$d\in\tfrac12\Z$, $0<d<w$, successive gaps are $2d$ and
$2(w-d)$, integers in $\{1,\ldots,2w-1\}$.
The residual shift in (5.6) allows no stronger universal margin;
strict inclusion and this finite set of gaps suffice.
The number $J$ of cuts is one of the even counts (5.4).
Each source lies wholly between consecutive operator cuts, because
the original clean cut and all its coarse images avoid sources.
At one cut the insertion boxes occupy disjoint spatial cells.

\section{Pressure, source comparison, and amplification}

\subsection{The only thermodynamic input}
Bounded finite-range Wilson interactions have a common pressure
$p_\beta$ on the coordinate and skew tori above:
\[
 \log Z_\Gamma[1]=p_\beta|\Gamma|+o(|\Gamma|).
 \tag{6.1}
\]
An elementary verification is to tile a fundamental domain by fixed
large coordinate boxes, deleting plaquettes at their boundaries.
Each deleted plaquette changes $\log Z$ by at most $\beta$.
Normalized Haar integrals factor on the disconnected boxes; leftover
links have integral one. First let the torus grow, so its boundary
fraction vanishes, then let the box size grow. The internal seam
fraction is $O(1/\text{box size})$. Comparing boxes also proves the
existence and uniqueness of their limiting pressure. The argument
works for bounded-aspect skew parallelepipeds and uniformly over
finitely many rectangular shapes. No quantitative convergence rate
or tilted pressure is needed.

\subsection{Full source arrays are essential}
Use the reflected source arrays $V_C,V_D$ in all $N$ fine cells.
Every fine-wall reflection is an integer affine signed permutation
preserving the coordinate Wilson law. Thus
$\mu_C(V_C)=N\mu_C(O)$, even though diagonal RP on $C_n$ has
not been assumed. Every cell factor $e^{tO^{[c]}}$ has the same
complete fold $e^{tV_D}$; (2.6) gives
\[
 Z_C[e^{tV_C}]\le Z_D[e^{tV_D}].
 \tag{6.2}
\]
Jensen and convexity of
$F_D(t)=\log(Z_D[e^{tV_D}]/Z_D[1])$ imply
\[
 \mu_C(O)\le\mu_{D,t}(O)
 +\frac{\log Z_D[1]-\log Z_C[1]}{tN},\qquad t>0.
 \tag{6.3}
\]
Here $F_D'(t)=N\mu_{D,t}(O)$ because the tilted law is invariant
under its wall group. Equation (6.1) makes the error tend to zero
at fixed $t$. A single-source unnormalized partition function raised
to $N$ would give an incorrect normalization; (6.2) uses full arrays
on both tori.

\subsection{Amplifying a clean-plane insertion}
Fix $G_X$ at one clean coordinate link plane and use the geometry
above. Coarse diagonal walls preserve $V_D$, so the normalized
coarse chessboard consequence of (2.6) is
\[
 |\mu_{D,t}(G_X)|^{N'}
 \le\frac{Z_D[e^{tV_D}\prod_gG_X^{[g]}]}{Z_D[e^{tV_D}]}.
 \tag{6.4}
\]
For clarity, apply (2.6) to one occupied coarse cell and the empty
ones: their complete folds contain respectively the $N'$ copies
and no insertion, while each contains the same source array.
This gives (6.4) with exactly the displayed denominator.

Apply the coordinate cell inequality to the numerator. If cell $c$
contains $m_c$ sources, its complete fold has $N''m_c$ source copies;
$\sum_cm_c=N$. If occupied, it also has $N''$ copies of its
operator, all of norm $\norm X$. The unnormalized bound (3.3)
therefore costs $e^{tN''m_c\norm O_\infty}\norm X^{N''}$
for occupied cells, and just the exponential for empty cells.
In that formula the source-free factors satisfy
\[
 \frac1J\sum_j\log Z_\Gamma(J\ell_j)
   =p_\beta|\Gamma|\sum_j\ell_j+o(n^4)
   =p_\beta n^4+o(n^4).
 \tag{6.5}
\]
Uniformity follows from the fixed finite set of gap lengths and
time directions. All spatial lengths and $J\ell_j$ grow in bounded
aspect ratios. After the $1/N''$ cell roots, the numerator bound is
\[
 \norm X^{N'}e^{tN\norm O_\infty}
       e^{p_\beta n^4+o(n^4)}.
\]
The one original denominator is at least
$e^{-tN\norm O_\infty}e^{p_\beta n^4+o(n^4)}$.
Taking the $N'$-th root proves
\[
 \limsup_{n\to\infty}|\mu_{D_n,t}(G_X)|
 \le e^{2t(H/h)^4\norm O_\infty}\norm X.
 \tag{6.6}
\]
All auxiliary sizes were fixed before this limit. In particular
$H$ may depend on the support of $X$, but not its norm. This
also proves (6.6) for the local unitization $I+isA$.

\section{Removing the tilt in one common state}

Fix a subsequence $\mu_{C_n}\to\nu$ and $t_j\downarrow0$.
For each $j$, extract a local weak limit $\lambda_j$ of
$\mu_{D_n,t_j}$ along a subsequence of those same volumes.
Then extract $\lambda_j\to\lambda$. Compactness follows from
the countable product of compact link groups. Equation (6.3)
gives $\nu(O)\le\lambda(O)$. All these laws are gauge invariant;
diagonal RP and invariance at the chosen fine wall survive both limits.

For every fixed finite kernel $X$, (6.6) holds along every volume
subsequence and yields
\[
 |\lambda(G_X)|\le\norm X.
 \tag{7.1}
\]
Thus one common $\lambda$ satisfies all these inequalities; no
operator-dependent choice of state is involved. Put
$\mathcal L(X)=\lambda(G_X)$, so $\mathcal L(I)=1$.
For self-adjoint $A$, apply (7.1) to $I+isA$:
\[
 |1+is\mathcal L(A)|^2\le1+s^2\norm A^2.
\]
Let $s\to0$ through both signs. Then
$\operatorname{Im}\mathcal L(A)=0$, and linearity gives
$\mathcal L(X^*)=\overline{\mathcal L(X)}$.
Reflection swaps the crossing kernel arguments, hence
$rG_X=\overline{G_{X^*}}$. Since $\lambda$ is a probability law,
\[
 \lambda(rG_X)=\lambda(G_X).
 \tag{7.2}
\]

Every smooth gauge-invariant two-layer observable $Y$ is uniformly
approximable by these insertions. After temporal parallel transport,
enlarge the selected spatial set so that $Y=Y(u,v)$, put
$K_Y(v,u)=K_S(v,u)Y(u,v)$, and let $P_m$ be product Peter--Weyl
cutoff. Its smooth kernels $P_mK_YP_m$ converge uniformly to $K_Y$.
The bounded finite-rank operator
\[
 X_m=C_S^{-1/2}P_mK_YP_mC_S^{-1/2}
\]
commutes with endpoint gauge transformations and has
$G_{X_m}=(P_mK_YP_m)/K_S\to Y$ uniformly. Positivity of every
retained convolution eigenvalue and of $K_S$ justifies both divisions.
No uniform bound on $\norm{X_m}$ is required: the tilt has already
been removed for each fixed operator. Uniform approximation extends
(7.2) to continuous $Y$.

Proposition~\ref{prop:layers}, followed by $t\downarrow0$ and
$\epsilon\downarrow0$, extends invariance to every local
gauge-invariant $P$ at the same clean plane. Gauge averaging a
cylinder test extends it to all local functions. Applying this
at two adjacent clean planes in each direction proves that
$\lambda$ is invariant under $2\Z^4$ translations. This step
establishes coordinate reflection \emph{invariance}; its diagonal
RP came from the adapted tilted measures.

\section{Convex separation proves the theorem}

Keep the desired mirror $x_1-x_2=0$ fixed. Let $\mathcal R$ be the
set of gauge-invariant probability laws invariant and RP at that
mirror, and invariant under $2\Z^4$. It is compact and convex:
on the compact configuration space, each invariance condition is a
closed linear equality and each continuous local RP test is a closed
affine inequality.

Given any real gauge-invariant local Peter--Weyl polynomial $P$,
translate it by an even vector into the positive side of the fixed
mirror, choose $h$ sufficiently large, and choose the other offsets
so its full support lies strictly inside that diagonal cell.
Even lattice translations form a full-rank sublattice, so this has
no parity obstruction. The source construction provides a law
$\lambda\in\mathcal R$ dominating the translated expectation.
Use the original law's full translation invariance and $\lambda$'s
even translation invariance to return to $P$. Therefore
\[
 \nu(P)\le\sup_{\lambda\in\mathcal R}\lambda(P)
 \quad\text{for every such }P.
 \tag{8.1}
\]

If $\nu\notin\mathcal R$, separation of the compact convex set
would give a strictly separating continuous real function. Cylinder
approximation, finite Peter--Weyl approximation, and then averaging
over the finite endpoint gauge group preserve strict separation and
produce a gauge-invariant finite polynomial.
All expectations under the gauge-invariant laws are unchanged by
that averaging. This contradicts (8.1). Hence $\nu\in\mathcal R$.
The coordinate tori already give $\nu$ full lattice translations,
signed permutations and coordinate RP. These transfer the proved
diagonal statement to the other mirrors and finish
Theorem~\ref{thm:main}.

This is averaging of a \emph{test function} in a separation argument,
not an assumption that translating or averaging RP states preserves
RP at a fixed plane. No uniqueness of a Gibbs state is used.

\section{What this buys for the continuum problem}

The result supplies a way to retain both coordinate and diagonal RP
in the same volume-first lattice theory. It repairs three concrete
interfaces: local gauge descent, a normalized clean-plane operator
bound, and the incompatible parity in the convenient QED3 placement.
The first comparison also requires full reflected source arrays;
omitting them loses the volume divisor in the pressure error.

For renormalized continuum observables, passing these inequalities
requires their common convergence on the relevant reflected test
algebra. That convergence and its all-order bounds remain dependent
on the model-specific RG estimates. Constants in the local reduction
may grow with $\beta$; this proof gives no uniform ultraviolet estimate.

The remaining main obligations are:
\begin{enumerate}
\item Verify the selected global four-layer complex-source chain and
the upstream running-coupling/cutoff input in the deposited account.
\item Prove the interacting angular analyticity and growth estimates
that would turn the two reflection directions into full rotations.
Notebook 01 explains why the reflection directions alone do not do so.
\item Establish an interacting nonzero positive-energy witness, then
construct and identify an adequate family of mixed gauge-invariant
observables. A single scalar sector is insufficient to identify
Yang--Mills theory.
\item Justify vacuum-sector reconstruction and separately prove any
claimed mass gap.
\end{enumerate}
No item in this list is discharged by the fixed-coupling theorem.

\begin{thebibliography}{9}
\bibitem{repo}
\emph{Yang--Mills LLM experiment}, public repository at
commit \texttt{e392ef4501\allowbreak 3a581664ba\allowbreak b76b908671\allowbreak 70b30fe10e}.
\url{https://github.com/ethanwsussman/Yang--Mills-LLM-experiment/tree/e392ef45013a581664bab76b90867170b30fe10e}.
\bibitem{qed}
\emph{QED3 account}, same repository,
\texttt{QED3\_slop/qed3.pdf}, Sections 3.2--3.3, printed pp.~20--25,
especially (3.17)--(3.27); Section 4.4, printed pp.~38--40,
especially (4.17)--(4.22). Source of the slab/comparison architecture;
its fermionic theorem is not imported as a Wilson proof.
\url{https://github.com/ethanwsussman/Yang--Mills-LLM-experiment/blob/e392ef45013a581664bab76b90867170b30fe10e/QED3_slop/qed3.pdf}.
\bibitem{clay}
A. Jaffe and E. Witten, \emph{Quantum Yang--Mills Theory}.
\url{https://www.claymath.org/wp-content/uploads/2022/06/yangmills.pdf}.
\end{thebibliography}
\end{document}
